\subsection{Factoring Numbers Review} 
Let's review the notion of factoring integers.
\begin{definition}
\begin{itemize*}
\item A \term{factor} (noun) of an integer $n$ is an integer $k$ such that $n=kd$ for some integer $d$.
\item To \term{factor} (verb) an integer is to determine its factors.
\item A \term{prime number} $p$ is an integer whose only factors are $\pm p$ or $\pm 1$.
\item The \term{prime factorization} of an integer as a product of prime numbers.
\item We can find the prime factorization of an integer by making a \term{factor tree}.
\end{itemize*}
\end{definition}
\begin{example}Make a factor tree to find the prime factorization of the number $18$.
\begin{lpic}{}{5}
\end{lpic}
\end{example}
\begin{theorem}[Fundamental Theorem of Arithmetic]
Any integer can be written uniquely as a product of prime numbers.
\end{theorem}
\begin{note}
For purposes of ``uniqueness,'' we ignore negative signs and order; so\\$10=2\cdot 5=5\cdot 2=(-2)\cdot(-5)$ would be considered ``the same'' factorization.
\end{note}